\section{Application}\label{appendix:application}

This appendix collects two constructions used in Section~\ref{sec:applications}: a smooth compact state space containing the relevant price box and a Banach structure on the parametrized class of dynamics.

\subsection{A smooth compact price domain}\label{appendix:application-smooth-domain}

Let $0<m<M<\infty$ be lower and upper bounds such that the equilibrium prices lie in $[m,M]^d\subset(0,\infty)^d$. Such bounds are supplied by \citet{skerlos2010fixed} for the mixed-logit pricing problem considered in Section~\ref{sec:applications}. We construct a compact smooth superset of this box that remains inside the positive orthant.

Fix an integer $p\ge 1$, let $c=(m+M)/2$, and define
\[
r=\frac{M-m}{2}\,d^{1/(2p)}.
\]
Set
\[
X_p=\left\{\vx\in\R^d:\;\sum_{i=1}^d \left(\frac{\vx_i-c}{r}\right)^{2p}\le 1\right\}.
\]
Then $X_p$ is compact and convex, and $[m,M]^d\subseteq X_p$. Indeed, if $\vx\in[m,M]^d$, then $|\vx_i-c|\le (M-m)/2$ for every $i$, so
\[
\begin{aligned}
\sum_{i=1}^d \left(\frac{\vx_i-c}{r}\right)^{2p}
&\le d\left(\frac{(M-m)/2}{r}\right)^{2p}\\
&=1.
\end{aligned}
\]

Moreover, $X_p$ has $C^\infty$ boundary. Its boundary is the level set of
\[
h(\vx)=\sum_{i=1}^d \left(\frac{\vx_i-c}{r}\right)^{2p},
\]
and $\nabla h(\vx)\ne 0$ on $h^{-1}(1)$ because $h(\vx)=1$ implies at least one coordinate satisfies $\vx_i\ne c$. Finally, for large enough $p$, $c-r>0$, since $d^{1/(2p)}\to 1$ as $p\to\infty$ and $c>(M-m)/2$. Hence $X_p\subset(0,\infty)^d$ for such $p$.

\subsection{A transported Banach structure}\label{appendix:application-banach}

Let
\[
\mathcal{Q}
=\left\{\vQ_{(\vtheta,\vomega)}:(\vtheta,\vomega)\in
\R^{d_C}\times\R^{d_\zeta}\right\}
\]
be the class of dynamics generated by the chosen parametrization of marginal costs and markups,
\[
\vQ_{(\vtheta,\vomega)}(\vx)
=\vc_{\vtheta}+\vzeta_{\vomega}(\vx)-\vx.
\]
Assume the parametrization $(\vtheta,\vomega)\mapsto \vQ_{(\vtheta,\vomega)}$ is injective. Transport the vector-space operations from $\R^{d_C}\times\R^{d_\zeta}$ to $\mathcal{Q}$ by defining
\[
\begin{aligned}
\vQ_{(\vtheta_1,\vomega_1)}\oplus \vQ_{(\vtheta_2,\vomega_2)}
&:=\vQ_{(\vtheta_1+\vtheta_2,\vomega_1+\vomega_2)},\\
\lambda\odot \vQ_{(\vtheta,\vomega)}
&:=\vQ_{(\lambda\vtheta,\lambda\vomega)}.
\end{aligned}
\]
Equip $\mathcal{Q}$ with the transported norm
\[
\|\vQ_{(\vtheta,\vomega)}\|_{\mathcal Q}
=\|(\vtheta,\vomega)\|_2.
\]
The injectivity assumption makes these operations and this norm well defined. With this structure, the parametrization is a linear isometry from $\R^{d_C}\times\R^{d_\zeta}$ onto $\mathcal{Q}$. Since $\R^{d_C}\times\R^{d_\zeta}$ is finite-dimensional and complete, $\mathcal{Q}$ is a Banach space.
